% ============================================
% 第5章：排名、积分、点赞：我们在设计"微型制度"
% 拱桥 · 激励设计
% ============================================

\chapter{排名、积分、点赞：我们在设计``微型制度''}

\vspace{0.5cm}

\begin{center}
{\small\color{muted} 拱桥的每一块石头都被两边的石头挤住。}
{\small\color{muted} 如果有一块石头想偷懒——它就会掉下来。}
{\small\color{muted} 制度，就是那两边的石头。}
\end{center}

\vspace{1cm}

\section{``分已经打完了，而且都是最高分''}

2026年6月23日，我在跟两位老师讨论我开发的AI教学平台。我讲了一段话，后来被转写成了文字：

\begin{shiyu-inline}
``以前我们在学习通里评分的时候，台上学生往台上一站，你不是可以评分了吗？分已经打完了，而且都是最高分。为啥他打多少分跟自己没关系？你再有一个是你只需要打个分嘛，写个分最省事，人就最小用力原则嘛。那只要打个分，打个满分，啪提交就完事了。''

\hfill ——2026年6月23日，AI教学平台讨论，录音转写
\end{shiyu-inline}

这段话描述的场景，每一个用过学习通评分功能的老师都见过。一个学生上台汇报，台下三十个学生，三十秒内全部打完分。清一色的满分。没有人真的在``评价''。所有人都在``完成''。

为什么？因为\textbf{打多少分，跟自己没关系。}

这就是旧系统里评价机制的致命缺陷：评价者和被评价者之间，没有利益关联。我给你打满分，不花我一分钱。我给你打低分，你也不会知道是我打的。那为什么不打满分呢？最小用力原则。

\section{``他必须打这个分很公正''}

然后我讲了我怎么改的：

\begin{shiyu-inline}
``但是现在我要做的是，因为你给别人打完分之后，别人马上加到积分里，这个积分一变就影响到了你自己的排名，影响了你期末的分数。所以他必须打这个分很公正。''

\hfill ——2026年6月23日，同一场讨论
\end{shiyu-inline}

我改了一个东西：\textbf{把评价者的利益和被评价者的表现绑在一起。}

你给别人打的分数，会影响别人的排名。别人的排名，会影响你的排名。如果你给一个很烂的作品打满分，那个人的排名会上升，你的排名会下降。如果你给一个很好的作品打低分，那个人的排名会下降——但你的排名不会因此上升，因为AI会检测你的评分是不是偏离了均值太远。

\textbf{这不是一个评分系统。这是一个微型制度。}

在这个制度里，``公正''不是一种道德要求，而是一种\textbf{利益计算}。你不需要做一个好人。你只需要做一个理性的人。而理性的人，会打公正的分。

\section{``AI给整理的所有学生对他的评价''}

制度里还有另一个细节：

\begin{shiyu-inline}
``30个评审团60句话。别人看到不是60句话，而是匿名的由AI整理出来的一段。AI给整理的所有学生对他的评价，整理完之后你看不出是哪位学生给他的这个东西，但一定里边你所说的话被采纳进去了。''

\hfill ——2026年6月23日，同一场讨论
\end{shiyu-inline}

我做了两件事：\textbf{匿名化}和\textbf{AI聚合}。

匿名化，让学生敢说真话。没有人知道那句批评是谁写的，所以没有人会因为批评同学而被记恨。

AI聚合，让学生的话被``听见''。三十个人写的六十句话，对一个人来说太多了。但AI把它们整理成一段流畅的反馈，那个学生就能在三十秒内读完，然后马上知道自己的问题在哪，下一版怎么改。

\textbf{制度的目的是让反馈变得可操作。}

\section{``他排名啪一杆就冲上去了''}

然后我讲了另一个设计：抢答机制。

\begin{shiyu-inline}
``第一个答完的学生能得十分，第二个答完能得九分，第三个答完得八分。咱以前课堂上最头疼的是好学生答完之后把答案啪一散，全班都知道。但是你现在没问题，你为了能够快速答完，好学生根本没工夫去管其他人。他赶紧先把自己答案提交了。''

``我上课经常有那样的学生，就是因为这一轮他答的最快最好，直接他就排名啪一杆就冲上去了。就这样，让有可能性他能翻盘的。所以他才愿意参与。要不然他发现自己反正我也超不过前面了，无所谓了，我就是躺平，我不跟你参与了。''

\hfill ——2026年6月23日，同一场讨论
\end{shiyu-inline}

这里有三个设计原则：

\textbf{第一，速度本身就是奖励。}第一个答对的人拿10分，第二个拿9分，第三个8分——后面的人拿到的分越来越低。这个设计把``快''变成了一种优势。好学生不会再把答案分享给全班——因为分享会降低自己的排名。

\textbf{第二，翻盘必须可能。}如果排名永远不变，排在后面的人就会放弃。但一轮抢答就能让排名``啪一杆就冲上去''——这个可能性，是让所有人持续参与的动力。

\textbf{第三，参与本身就有分。}不是只有答对才得分。投票、开放式问题、甚至点``不懂''——只要参与了，AI就给你分。因为\textbf{参与本身就是价值}。一个沉默的课堂，比一个答错的课堂更可怕。

\begin{insight}
\label{insight:chap5}

\textbf{积分系统不是在管理学生。积分系统是在设计课堂的微观制度。}

每一项规则——谁得分、怎么得分、得分之后会发生什么——都在塑造学生的行为。如果你奖励``正确''，学生就会追求安全。如果你奖励``快速''，学生就会追求效率。如果你奖励``参与''，学生就会追求表达。

制度不是中性的。制度本身就在教东西。
\end{insight}

\section{制度的双面性}

2026年6月，在教学理念综合分析中，有一段话让我反复读：

\begin{shiyu-inline}
``积分系统实际上是在设计课堂的'微型制度'。有的平台把学生给他人的评分与自己的排名、期末成绩关联，解决'人人随手满分'；点赞、评论、与教师判断一致还会产生不同积分。这非常聪明，因为它改变了参与的激励结构；也非常危险，因为学生可能开始优化积分、迎合教师或交换人情，而非追求真实判断。''

\hfill ——教学理念综合分析，2026年7月
\end{shiyu-inline}

\textbf{制度是双面的。}

它能让沉默的课堂变活跃，也能让活跃的课堂变表演。它能让学生从不参与到参与，也能让学生从``参与''到``迎合''。它能解决``人人随手满分''，也可能制造``人人优化积分''。

这不是制度的错。任何制度都有双面性。问题不是``要不要制度''，而是\textbf{制度有没有被有意识地设计，以及有没有被持续地审视}。

在我的课上，我保留了那个反思：\textbf{即使平台高度智能化，``人与人的东西''仍被认为最重要。}积分是辅助，不是目的。学生迭代的过程应该由学生自己总结，而不是外包给AI。

\section{拱桥的隐喻之二}

拱桥的每一块石头，都被两边的石头挤住。如果有一块石头想偷懒——它就会掉下来。制度，就是那两边的石头。

但制度不是越多越好。石头太密，拱桥就变成了实心的墙——没有空间，没有呼吸，没有光。制度太松，拱桥就散了——每一块石头都只为自己，整座桥就塌了。

\textbf{好的制度，是让每个人在追求自身利益的同时，恰好做了对整体有利的事。}

就像拱桥。每一块石头都在承受自己的重量。但因为它被两边的石头挤住了，它承受的重量，恰好帮助了整座桥站住。

\begin{shiyu-note}
\label{note:chap5}

\textbf{制度是双面的。但逃避制度的设计，是更危险的。}

很多老师不愿意设计积分系统。他们觉得``学习应该是自发的''，``积分会让学习变味''。但问题是：\textbf{不做设计的课堂，也有制度。}它只是隐性的——学生知道谁在认真听、谁在刷手机、谁会被老师喜欢、谁会被忽视。这些隐性的制度，比显性的积分系统更不公平，因为没有人能质疑它。

把制度显性化，才是公平的开始。你可以看到规则，你可以质疑规则，你可以建议修改规则。制度变成了一件可以被讨论和迭代的东西，而不是一件``从来就是这样''的天经地义。
\end{shiyu-note}

\vspace{1cm}

\begin{decision-point}
\label{decision:chap5}

\noindent\textbf{这一章，我们看到了积分系统如何设计课堂的微型制度。}

\vspace{0.3cm}

\noindent 下一次你让学生互评的时候，不要只让他们打分。加一句话：``你认为这个作品最大的优点是什么？如果你来改，你会改哪里？''

\noindent 然后，把这些评语匿名汇总，用AI整理成一段话，发给那个学生。告诉他：这是三十个人帮你想的改进方向。

\noindent 看看他下一版会改成什么样。
\end{decision-point}

\vspace{0.5cm}

\begin{flushright}
{\small\color{muted} 下一章：美丽的初稿——AI把``做出来''的门槛降为零，把``知道好不好''抬向无限。}
\end{flushright}